\documentclass[10pt,a4paper]{amsart}
\usepackage[margin=19mm]{geometry}
\usepackage{amsmath,amssymb,mathtools}
\usepackage[scr=rsfs,cal=euler]{mathalfa}
\usepackage{xfrac}
\usepackage[unicode,hidelinks,bookmarks=false]{hyperref}
\newcommand{\ie}{i.e.\ }
\newcommand{\cf}{cf.\ }
\newcommand{\resp}{respectively\ }
\newcommand{\Subsection}{\subsection}
\providecommand{\nobreakdash}{\nobreak\hbox{-}}
\setlength{\emergencystretch}{2em}
\multlinegap0pt
\usepackage{bm}
\usepackage{bbm}
\usepackage{dsfont}
\usepackage{xspace}
\usepackage{cancel}
\usepackage{mathtools}
\usepackage{amsthm}
\makeatletter
\@ifundefined{theorem}{\newtheorem{theorem}{Theorem}}{}
\@ifundefined{assumption}{\newtheorem{assumption}[theorem]{Assumption}}{}
\@ifundefined{lemma}{\newtheorem{lemma}[theorem]{Lemma}}{}
\makeatother
\newcommand{\EE}{\mathbb{E}}
\newcommand{\RD}{\mathbb{R}^{d}}
\title[Plug-In Classification of Drift Functions in Diffusion Proce]{Plug-In Classification of Drift Functions in Diffusion Processes Using Neural Networks}
\author{Yuzhen Zhao, Jiarong Fan, Yating Liu}
\date{Source: arXiv:2602.02791v2 (CC BY 4.0)}
\begin{document}
\maketitle

\noindent\emph{Source and licence.} Plug-In Classification of Drift Functions in Diffusion Processes Using Neural Networks. Version arXiv:2602.02791v2 (CC BY 4.0), distributed under the Creative Commons Attribution 4.0 International licence, as recorded in the arXiv licence metadata for this exact version and shown on the abstract page https://arxiv.org/abs/2602.02791v2. Full rights notice: licenses/arxiv-2602.02791-CC-BY-4.0.txt.\par\smallskip

\noindent\textit{Editorial scope.} Counted unit: the Lemma whose source label is \texttt{lem:a-hat-sigma-hat-bound} in the article (source0.tex lines 2086--2095, source label \texttt{lem:a-hat-sigma-hat-bound}), delivered together with the article's own complete original proof of it (source0.tex lines 2097--2154, one \texttt{proof} environment of 58 lines). No mathematics is written, repaired or re-derived: statement and proof are the article's own text line for line. Required context retained uncounted: assump:main-theorem (lines 443--445); assump:sigma-hat (lines 1974--1979). Every definition, display and label of those lines is retained unchanged. Omitted from this package and not redistributed: 1--442, 446--1973, 1980--2085, 2096--2096, 2155--2843. Nothing inside a retained excerpt is omitted or altered: every retained body is delivered exactly as the article's lines are written, so no omission receipt of any kind accompanies this package. Citations: the retained excerpts carry no citation command at all, so no bibliography entry of the article is transcribed and no reference is redistributed. Reference adaptation: none was needed and none was made; every reference inside the delivered unit resolves inside it, so no unresolved reference or citation is produced. Printed slips kept literal: no printed word, symbol, hypothesis or number of the counted statement or of its proof was altered. Layout and class: the article's own class is \texttt{article} with packages including neurips\_2026, inputenc, fontenc, hyperref, url, booktabs, amsfonts, nicefrac, microtype, xcolor, graphicx, subcaption, amsmath, amssymb; the wrapper is the lane's pinned \texttt{amsart} editorial wrapper at 10pt on A4, which restates the theorem-like environments the retained text needs and adds the article's own macro definitions that the retained text needs (\texttt{\textbackslash EE}, \texttt{\textbackslash RD}), with extra wrapper packages (bm, bbm, dsfont, xspace, cancel, mathtools). The delivered numbering is the unit's own. Assets: no retained span contains a figure environment, a graphics-inclusion command, a PDF-page inclusion, a table or a TikZ picture; no image file is copied into this package, the package declares no asset dependency and no image file is redistributed with it. Licence: version arXiv:2602.02791v2 of this article is distributed under the Creative Commons Attribution 4.0 International licence, as recorded in the arXiv licence metadata for this exact version and shown on the abstract page https://arxiv.org/abs/2602.02791v2. A full rights notice is delivered at licenses/arxiv-2602.02791-CC-BY-4.0.txt. Characters: every retained excerpt is pure ASCII, so the delivered engine, XeLaTeX with its default Latin Modern text font, reports no missing character.
\begin{assumption}\label{assump:main-theorem}
For every $x \in \mathbb{R}^d$, the matrix $a(x) \coloneqq \sigma\sigma^{\top}(x)$ is invertible, and $a^{-1}$ is globally Lipschitz continuous with respect to the operator norm $|\cdot|_{\mathrm{op}}$, with Lipschitz constant $L_{a^{-1}}$. Moreover, there exists a constant $\Lambda > 0$ such that $\max(|\sigma(x)|,|a(x)^{-1}|_{\mathrm{op}}) \leq \Lambda, \,x\in\RD$.
\end{assumption}

\begin{assumption}\label{assump:sigma-hat}
The estimator $\widehat\sigma(x)$ is invertible for every $x\in\RD$ and satisfies
\[
\forall x\in\RD,\qquad \max(|\widehat\sigma(x)|, |\widehat{a}(x)^{-1}|_{\mathrm{op}}) \le \Lambda.
\]
\end{assumption}

\begin{lemma}\label{lem:a-hat-sigma-hat-bound}
Assume that Assumptions~\ref{assump:main-theorem} and \ref{assump:sigma-hat} hold. We have
\[
\EE\!\left[
\frac{1}{M}\sum_{m=0}^{M-1}
\left|\widehat a(X_{t_m})-a(X_{t_m})\right|_{\mathrm{op}}^2
\right]
\le 4\Lambda^2\,\mathcal E(\widehat \sigma,\sigma).
\]
\end{lemma}

\begin{proof}[Proof of Lemma \ref{lem:a-hat-sigma-hat-bound}]
For every \(x\in\RD\), we write
\begin{align*}
\widehat a(x)-a(x)
&=
\widehat \sigma(x)\widehat \sigma(x)^{\top}
-\sigma(x)\sigma(x)^{\top}=
\bigl(\widehat \sigma(x)-\sigma(x)\bigr)\widehat \sigma(x)^{\top}
+\sigma(x)\bigl(\widehat \sigma(x)-\sigma(x)\bigr)^{\top}.
\end{align*}
Taking the operator norm and using the triangle inequality, together with the submultiplicativity of the operator norm, yields
\begin{align*}
\left|\widehat a(x)-a(x)\right|_{\mathrm{op}}
&\le
\left|\bigl(\widehat \sigma(x)-\sigma(x)\bigr)\widehat \sigma(x)^{\top}\right|_{\mathrm{op}}
+
\left|\sigma(x)\bigl(\widehat \sigma(x)-\sigma(x)\bigr)^{\top}\right|_{\mathrm{op}}\\
&\le
\left|\widehat \sigma(x)-\sigma(x)\right|_{\mathrm{op}}
\left|\widehat \sigma(x)^{\top}\right|_{\mathrm{op}}
+
\left|\sigma(x)\right|_{\mathrm{op}}
\left|\bigl(\widehat \sigma(x)-\sigma(x)\bigr)^{\top}\right|_{\mathrm{op}}.
\end{align*}
%Since the operator norm is invariant under transposition, we have
%\[
%\left|\widehat \sigma(x)^{\top}\right|_{\mathrm{op}}
%=
%\left|\widehat \sigma(x)\right|_{\mathrm{op}},
%\qquad
%\left|\bigl(\widehat \sigma(x)-\sigma(x)\bigr)^{\top}\right|_{\mathrm{op}}
%=
%\left|\widehat \sigma(x)-\sigma(x)\right|_{\mathrm{op}}.
%\]
Therefore,
\[
\left|\widehat a(x)-a(x)\right|_{\mathrm{op}}
\le
\bigl(
|\widehat \sigma(x)|_{\mathrm{op}}+|\sigma(x)|_{\mathrm{op}}
\bigr)
\left|\widehat \sigma(x)-\sigma(x)\right|_{\mathrm{op}}\le 2\Lambda \left|\widehat \sigma(x)-\sigma(x)\right|_{\mathrm{op}}
\]
and
\[
\EE\!\left[
\frac{1}{M}\sum_{m=0}^{M-1}
\left|\widehat a(X_{t_m})-a(X_{t_m})\right|_{\mathrm{op}}^2
\right]
\le
4\Lambda^2
\EE\!\left[
\frac{1}{M}\sum_{m=0}^{M-1}
\left|\widehat \sigma(X_{t_m})-\sigma(X_{t_m})\right|_{\mathrm{op}}^2
\right]= 4\Lambda^2\,\mathcal E(\widehat \sigma,\sigma). \hfill\qedhere
\]
%The last claim follows immediately from the definition of \(\mathcal E(\widehat \sigma,\sigma)\).
\end{proof}

\end{document}
